\documentclass[10pt, aspectratio=169]{beamer}
\usepackage{docmute}
\usepackage{beamer_theme_408}

\title{408考研零基础 --- 数据结构}
\author{}
\date{}

\begin{document}

\frame{\titlepage}

\begin{frame}{模块1 数据结构}
    本模块包含以下 18 节内容：
    \begin{enumerate}
        \item 1-1 数据结构的基本概念
        \item 1-2 时间复杂度与空间复杂度
        \item 1-3 线性表概念与顺序存储
        \item 1-4 线性表的链式存储
        \item 1-5 链表进阶
        \item 1-6 栈
        \item 1-7 队列
        \item 1-8 数组与特殊矩阵
        \item 1-9 树的基本概念
        \item 1-10 二叉树
        \item 1-11 二叉树的遍历
        \item 1-12 线索二叉树
        \item 1-13 树、森林与二叉树
        \item 1-14 哈夫曼树与并查集
        \item 1-15 图的基本概念与存储
        \item 1-16 图的遍历
        \item 1-17 查找基础
        \item 1-18 排序入门
    \end{enumerate}
\end{frame}

\subsection{1-1 数据结构的基本概念}
\documentclass[10pt, aspectratio=169]{beamer}

\usetheme{Madrid}
\usecolortheme{dolphin}

\usepackage{ctex}
\usepackage{xeCJK}
\usepackage{amsmath, amssymb}
\usepackage{graphicx}
\usepackage{listings}
\usepackage{tikz}
\usepackage{xcolor}

\title{408考研零基础 --- 数据结构}
\subtitle{1-1 数据结构的基本概念}
\author{}
\date{}

\setbeamertemplate{page number in head/foot}{}

\begin{document}


\begin{frame}{本节目标}
    \begin{enumerate}
        \item 掌握数据、数据元素、数据项的定义及其层级关系
        \item 理解数据结构的二元组形式化定义
        \item 区分逻辑结构与存储结构，掌握四种逻辑结构与两种存储结构的特征
        \item 理解抽象数据类型 ADT 的概念
    \end{enumerate}
\end{frame}

\begin{frame}{数据的定义}
    \textbf{数据（Data）}
    \begin{itemize}
        \item 客观事物的符号表示
        \item 所有能输入到计算机中并被计算机程序处理的符号的总称
    \end{itemize}
    \vspace{0.5em}
    \textbf{分类}
    \begin{itemize}
        \item 数值数据：整数、实数
        \item 非数值数据：字符、图像、声音
    \end{itemize}
\end{frame}

\begin{frame}{数据元素与数据项}
    \textbf{数据元素（Data Element）}
    \begin{itemize}
        \item 数据的基本单位
        \item 在计算机程序中通常作为一个整体进行考虑和处理
        \item 一个数据元素可由若干个数据项组成
    \end{itemize}
    \vspace{0.5em}
    \textbf{数据项（Data Item）}
    \begin{itemize}
        \item 数据的不可分割的最小单位
    \end{itemize}
    \vspace{0.5em}
    \textbf{数据对象（Data Object）}
    \begin{itemize}
        \item 性质相同的数据元素的集合
        \item 是数据的一个子集
    \end{itemize}
\end{frame}

\begin{frame}{四概念层级关系}
    \begin{center}
    \begin{tikzpicture}
        \draw[blue, thick] (0,0) ellipse (3cm and 2cm);
        \node[blue] at (0,1.2) {数据 (Data)};
        \draw[blue, thick] (0,0) ellipse (2cm and 1.3cm);
        \node[blue] at (0,0.6) {数据对象 (Data Object)};
        \draw[blue, thick] (0,0) ellipse (1.2cm and 0.7cm);
        \node[blue] at (0,0) {数据元素};
        \draw[blue, thick] (0,0) ellipse (0.5cm and 0.3cm);
        \node[blue, font=\tiny] at (0,-0.2) {数据项};
    \end{tikzpicture}
    \end{center}
    \begin{center}
        \textcolor{blue}{数据 $\supset$ 数据对象 $\supset$ 数据元素 $\supset$ 数据项}
    \end{center}
\end{frame}

\begin{frame}{数据结构的定义}
    \textbf{形式化定义}
    \begin{itemize}
        \item 数据结构是一个二元组：
            \[
            \text{Data\_Structure} = (D, S)
            \]
        \item $D$ --- 数据元素的有限集合
        \item $S$ --- $D$ 上关系的有限集合
    \end{itemize}
    \vspace{0.5em}
    \textbf{关系的形式化表示}
    \[
    D = \{a_1, a_2, \ldots, a_n\}, \quad S = \{R\}
    \]
    \[
    r = \{\langle a_i, a_j \rangle \mid a_i, a_j \in D\}
    \]
\end{frame}

\begin{frame}{数据结构的两要素}
    \begin{center}
    \begin{tikzpicture}
        \draw[blue, thick] (0,0) rectangle (4,2);
        \node[blue] at (2,1.5) {\textbf{数据结构}};
        \draw[->, blue, thick] (0.5,0.5) -- (0.5,-0.5);
        \draw[->, blue, thick] (3.5,0.5) -- (3.5,-0.5);
        \node at (0.5,-0.8) {数据元素集合 $D$};
        \node at (3.5,-0.8) {关系集合 $S$};
        \node[font=\small] at (2,-1.5) {二者缺一不可};
    \end{tikzpicture}
    \end{center}
\end{frame}

\begin{frame}{逻辑结构 --- 分类}
    \textbf{逻辑结构}：描述数据元素之间的逻辑关系，与存储方式无关
    \begin{enumerate}
        \item \textbf{集合结构}：元素间除同属集合外无其他关系
        \item \textbf{线性结构}：一对一关系，唯一前驱/后继
        \item \textbf{树形结构}：一对多层次关系
        \item \textbf{图状结构}：多对多任意关系
    \end{enumerate}
\end{frame}

\begin{frame}{四种逻辑结构对比}
    \begin{center}
    \begin{tabular}{|c|c|c|}
        \hline
        \textbf{类型} & \textbf{关系特征} & \textbf{典型实例} \\
        \hline
        集合结构 & 无关系 & 同班学生 \\
        \hline
        线性结构 & 一对一 & 线性表、栈、队列 \\
        \hline
        树形结构 & 一对多 & 二叉树、B树 \\
        \hline
        图状结构 & 多对多 & 有向图、无向图 \\
        \hline
    \end{tabular}
    \end{center}
\end{frame}

\begin{frame}{存储结构 --- 顺序存储}
    \textbf{顺序存储结构}
    \begin{itemize}
        \item 借助元素在存储器中的相对位置表示逻辑关系
        \item 通常使用数组实现
        \item 特征：逻辑相邻 $\implies$ 物理相邻
    \end{itemize}
    \vspace{0.5em}
    \begin{center}
    \begin{tabular}{|c|c|}
        \hline
        \textbf{优点} & \textbf{缺点} \\
        \hline
        随机访问 O(1) & 插入删除 O(n) \\
        \hline
    \end{tabular}
    \end{center}
\end{frame}

\begin{frame}{存储结构 --- 链式存储}
    \textbf{链式存储结构}
    \begin{itemize}
        \item 借助指针表示元素之间的逻辑关系
        \item 每个结点包含数据域与指针域
        \item 特征：逻辑相邻 $\nRightarrow$ 物理相邻
    \end{itemize}
    \vspace{0.5em}
    \begin{center}
    \begin{tabular}{|c|c|}
        \hline
        \textbf{优点} & \textbf{缺点} \\
        \hline
        插入删除 O(1) & 不支持随机访问 \\
        \hline
    \end{tabular}
    \end{center}
\end{frame}

\begin{frame}{ADT --- 抽象数据类型}
    \textbf{定义}：一个数学模型以及定义在该模型上的一组操作
    \[
    \text{ADT} = (D, S, P)
    \]
    \begin{itemize}
        \item $D$ --- 数据对象
        \item $S$ --- $D$ 上的关系集合
        \item $P$ --- 对 $D$ 的基本操作集合
    \end{itemize}
    \vspace{0.5em}
    \textbf{核心特征}：
    \begin{enumerate}
        \item 数据抽象 --- 仅描述逻辑特性
        \item 信息隐藏 --- 实现细节对外不可见
        \item 接口与实现分离
    \end{enumerate}
\end{frame}

\begin{frame}{栈的 ADT 定义示例}
    \begin{center}
    \begin{minipage}{0.85\textwidth}
    \begin{verbatim}
ADT Stack {
    数据对象: D = {ai | ai in ElemSet}
    数据关系: R = {<ai-1, ai>}
    基本操作:
      InitStack(&S)    构造空栈
      Push(&S, e)      入栈
      Pop(&S, &e)      出栈
      GetTop(S, &e)    取栈顶元素
      StackEmpty(S)    判栈空
}
    \end{verbatim}
    \end{minipage}
    \end{center}
\end{frame}

\begin{frame}{例题 --- 逻辑结构判定}
    \textbf{题目}：判断下列数据集合的逻辑结构类型
    \begin{enumerate}
        \item 某班级全体学生的学号集合，元素间无特定关系
        \item 线性队列中等待处理的任务序列
        \item 文件系统中的目录与子目录关系
        \item 城市之间的交通路线图
    \end{enumerate}
    \vspace{0.5em}
    \textbf{答案}：
    \begin{enumerate}
        \item 集合结构
        \item 线性结构
        \item 树形结构
        \item 图状结构
    \end{enumerate}
\end{frame}

\begin{frame}{本节总结}
    \begin{enumerate}
        \item 数据结构定义为二元组 $(D, S)$
        \item 逻辑结构：集合、线性、树形、图状
        \item 存储结构：顺序存储、链式存储
        \item 逻辑结构与存储结构相互独立
        \item ADT 实现数据模型与操作的封装
    \end{enumerate}
\end{frame}

\end{document}


\subsection{1-2 时间复杂度与空间复杂度}
\documentclass[10pt, aspectratio=169]{beamer}
\usepackage{../beamer_theme_408}

\title{408考研零基础 --- 数据结构}
\subtitle{1-2 时间复杂度与空间复杂度}
\author{}
\date{}


\begin{document}


\begin{frame}{本节目标}
    \begin{enumerate}
        \item 理解算法与程序的区别，掌握算法的五个特征
        \item 理解时间复杂度的定义与大 O 记号的含义
        \item 能够分析常见算法的时间复杂度
        \item 理解空间复杂度的定义及其分析方法
    \end{enumerate}
\end{frame}

\begin{frame}{算法的基本概念}
    \textbf{算法}：对特定问题求解步骤的一种描述，是指令的有限序列
    \vspace{0.5em}
    \textbf{五个特征}
    \begin{enumerate}
        \item \textbf{有穷性}：必须在执行有穷步后结束
        \item \textbf{确定性}：每条指令有确切含义，相同输入有相同输出
        \item \textbf{可行性}：操作可以通过已实现的基本运算完成
        \item \textbf{输入}：零个或多个输入
        \item \textbf{输出}：一个或多个输出
    \end{enumerate}
    \vspace{0.5em}
    \begin{block}{注意}
        算法 $\neq$ 程序，程序不一定满足有穷性
    \end{block}
\end{frame}

\begin{frame}{时间复杂度的定义}
    \textbf{时间复杂度} $T(n)$
    \begin{itemize}
        \item 算法中基本运算的频度
        \item $n$ 为问题的规模
    \end{itemize}
    \vspace{0.5em}
    \textbf{大 O 记号}
    \[
    T(n) = O(f(n)) \iff \exists c > 0, n_0 > 0 \text{ s.t. } T(n) \le c\cdot f(n), \forall n \ge n_0
    \]
    \vspace{0.5em}
    \textbf{常见时间复杂度（递增）}
    \[
    O(1) < O(\log n) < O(n) < O(n\log n) < O(n^2) < O(2^n) < O(n!)
    \]
\end{frame}

\begin{frame}{时间复杂度分析 --- 案例}
    \textbf{案例一：顺序查找}
    \begin{itemize}
        \item 最好：$O(1)$
        \item 最坏：$O(n)$
        \item 平均：$O(n)$
    \end{itemize}
    \vspace{0.5em}
    \textbf{案例二：两层嵌套循环}
    \[
    \sum_{i=1}^{n}\sum_{j=1}^{n} 1 = n^2 \implies O(n^2)
    \]
    \vspace{0.5em}
    \textbf{案例三：二分查找}
    \[
    \frac{n}{2^k} = 1 \implies k = \log_2 n \implies O(\log n)
    \]
\end{frame}

\begin{frame}{空间复杂度}
    \textbf{定义}：算法运行时所需的存储空间大小
    \[
    S(n) = O(f(n))
    \]
    \vspace{0.5em}
    \textbf{组成部分}
    \begin{itemize}
        \item 输入数据所占空间
        \item 程序本身所占空间
        \item 辅助变量所占空间（关注重点）
    \end{itemize}
    \vspace{0.5em}
    \textbf{常见情况}
    \begin{itemize}
        \item 原地工作：$O(1)$
        \item 递归栈深度为 $n$：$O(n)$
    \end{itemize}
\end{frame}

\begin{frame}{例题分析}
    \textbf{例题一}
    \[
    \begin{aligned}
    &\text{for}(i = 1; i \le n; i++) \\
    &\qquad \text{for}(j = 1; j \le i; j++) \\
    &\qquad \qquad x\text{++};
    \end{aligned}
    \]
    基本运算次数：$1 + 2 + \cdots + n = n(n+1)/2$
    \[
    T(n) = O(n^2)
    \]
    \vspace{0.5em}
    \textbf{例题二：斐波那契递归}
    \[
    T(n) = T(n-1) + T(n-2) \implies T(n) = O(2^n)
    \]
\end{frame}

\begin{frame}{本节总结}
    \begin{enumerate}
        \item 算法特征：有穷性、确定性、可行性、输入、输出
        \item 大 O 表示时间复杂度的数量级
        \item 常见复杂度：$O(1) < O(\log n) < O(n) < O(n^2) < O(2^n)$
        \item 空间复杂度衡量辅助存储空间
        \item 分析步骤：确定问题规模 $\to$ 找出基本运算 $\to$ 建立函数关系 $\to$ 取数量级
    \end{enumerate}
\end{frame}

\end{document}


\subsection{1-3 线性表概念与顺序存储}
\documentclass[10pt, aspectratio=169]{beamer}

\usetheme{Madrid}
\usecolortheme{dolphin}

\usepackage{ctex}
\usepackage{xeCJK}
\usepackage{amsmath, amssymb}
\usepackage{graphicx}
\usepackage{listings}
\usepackage{tikz}
\usepackage{xcolor}

\title{408考研零基础 --- 数据结构}
\subtitle{1-3 线性表概念与顺序存储}
\author{}
\date{}

\setbeamertemplate{page number in head/foot}{}

\begin{document}

\frame{\titlepage}

\begin{frame}{本节目标}
    \begin{enumerate}
        \item 掌握线性表的形式化定义与逻辑特征
        \item 理解顺序表的存储结构及其优缺点
        \item 掌握顺序表的插入、删除、查找操作的实现与时间复杂度分析
    \end{enumerate}
\end{frame}

\begin{frame}{线性表的定义}
    \textbf{线性表}：具有相同数据类型的 $n$ 个数据元素的有限序列
    \[
    L = (a_1, a_2, \ldots, a_i, \ldots, a_n)
    \]
    \vspace{0.5em}
    \textbf{逻辑特征}
    \begin{enumerate}
        \item 存在唯一的第一个元素 $a_1$（表头）
        \item 存在唯一的最后一个元素 $a_n$（表尾）
        \item 除第一个元素外，每个元素有且仅有一个直接前驱
        \item 除最后一个元素外，每个元素有且仅有一个直接后继
    \end{enumerate}
    \vspace{0.5em}
    $n$ 为表长，$n = 0$ 时称为空表
\end{frame}

\begin{frame}{线性表的 ADT 定义}
    \begin{block}{ADT List}
        数据对象：$D = \{a_i \mid a_i \in \text{ElemSet}, i = 1,2,\ldots,n, n \ge 0\}$ \\
        数据关系：$R = \{\langle a_{i-1}, a_i \rangle \mid a_{i-1}, a_i \in D, i = 2,3,\ldots,n\}$
    \end{block}
    \vspace{0.5em}
    \textbf{基本操作}
    \begin{itemize}
        \item InitList(\&L) --- 初始化
        \item ListInsert(\&L, i, e) --- 插入
        \item ListDelete(\&L, i, \&e) --- 删除
        \item LocateElem(L, e) --- 按值查找
        \item GetElem(L, i) --- 按位查找
    \end{itemize}
\end{frame}

\begin{frame}{顺序表的结构定义}
    \textbf{顺序表}：用一组地址连续的存储单元依次存储线性表的数据元素
    \[
    \text{LOC}(a_i) = \text{LOC}(a_1) + (i-1) \times k
    \]
    \vspace{0.5em}
    \textbf{C 语言定义}
    \begin{lstlisting}[language=C, basicstyle=\ttfamily\small]
#define MaxSize 50
typedef struct {
    ElemType data[MaxSize];
    int length;
} SqList;
    \end{lstlisting}
    \vspace{0.5em}
    \begin{block}{核心特性}
        随机访问 --- 时间复杂度 $O(1)$
    \end{block}
\end{frame}

\begin{frame}{顺序表 --- 插入操作}
    \textbf{在第 $i$ 个位置插入新元素 $e$}
    \vspace{0.3em}
    \textbf{步骤}
    \begin{enumerate}
        \item 判断 $i$ 是否合法（$1 \le i \le \text{length}+1$）
        \item 判断表是否已满
        \item 将第 $i$ 至第 $n$ 个元素依次后移
        \item 插入 $e$，表长加一
    \end{enumerate}
    \vspace{0.5em}
    \textbf{时间复杂度}
    \[
    \text{平均移动次数} = \frac{n}{2} \implies O(n)
    \]
\end{frame}

\begin{frame}{顺序表 --- 删除操作}
    \textbf{删除第 $i$ 个位置的元素}
    \vspace{0.3em}
    \textbf{步骤}
    \begin{enumerate}
        \item 判断 $i$ 是否合法（$1 \le i \le \text{length}$）
        \item 将第 $i+1$ 至第 $n$ 个元素依次前移
        \item 表长减一
    \end{enumerate}
    \vspace{0.5em}
    \textbf{时间复杂度}
    \[
    \text{平均移动次数} = \frac{n-1}{2} \implies O(n)
    \]
\end{frame}

\begin{frame}{顺序表 --- 按值查找}
    \textbf{查找第一个值为 $e$ 的元素并返回位序}
    \vspace{0.3em}
    从第一个元素开始依次比较
    \vspace{0.5em}
    \textbf{时间复杂度}
    \begin{itemize}
        \item 最好情况：$O(1)$
        \item 最坏情况：$O(n)$
        \item 平均情况：$O(n)$
    \end{itemize}
    \vspace{0.5em}
    \begin{block}{顺序表总结}
        \begin{tabular}{ll}
            优点：随机访问 $O(1)$，存储密度高 \\
            缺点：插入删除 $O(n)$，需预分配空间
        \end{tabular}
    \end{block}
\end{frame}

\begin{frame}{例题 --- 删除所有值为 x 的元素}
    \textbf{设计算法，删除顺序表中所有值为 $x$ 的元素}
    \[
    O(n), O(1)
    \]
    \vspace{0.3em}
    \textbf{思路}：覆盖代替多次移动
    \vspace{0.3em}
    \begin{enumerate}
        \item 设置指针 $k = 0$
        \item 遍历顺序表
        \item 若 $A[i] \ne x$，则 $A[k] = A[i], k++$
        \item 跳过等于 $x$ 的元素
        \item 表长设为 $k$
    \end{enumerate}
\end{frame}

\begin{frame}{本节总结}
    \begin{enumerate}
        \item 线性表是 $n$ 个相同类型数据元素的有限序列，一对一关系
        \item 顺序表采用连续存储，支持 $O(1)$ 随机访问
        \item 插入和删除平均移动 $n/2$ 个元素，时间复杂度 $O(n)$
        \item 按值查找平均 $O(n)$
        \item 顺序表适用于查找频繁、插入删除较少的场景
    \end{enumerate}
    \vspace{1em}
    \begin{center}
        \textcolor{blue}{\textbf{下节预告：线性表的链式存储}}
    \end{center}
\end{frame}

\end{document}


\subsection{1-4 线性表的链式存储}
\documentclass[10pt, aspectratio=169]{beamer}
\usepackage{../beamer_theme_408}

\title{408考研零基础 --- 数据结构}
\subtitle{1-4 线性表的链式存储}
\author{}
\date{}


\begin{document}


\begin{frame}{本节目标}
    \begin{enumerate}
        \item 理解单链表的结点结构，掌握头指针与头结点的概念
        \item 掌握单链表的建立方法（头插法和尾插法）
        \item 掌握单链表的插入、删除、查找操作的实现与时间复杂度分析
    \end{enumerate}
\end{frame}

\begin{frame}[fragile]{单链表的结构定义}
    每个结点包含数据域和指针域
    \vspace{0.5em}
    \begin{lstlisting}[language=C, basicstyle=\ttfamily\small]
typedef struct LNode {
    ElemType data;
    struct LNode *next;
} LNode, *LinkList;
    \end{lstlisting}
    \vspace{0.5em}
    \textbf{重要概念}
    \begin{itemize}
        \item \textbf{头指针}：指向链表第一个结点的指针
        \item \textbf{头结点}：第一个结点前的附加结点
    \end{itemize}
    \vspace{0.3em}
    \begin{block}{引入头结点的好处}
        第一个元素前插入或删除第一个元素时，操作与其他位置一致
    \end{block}
\end{frame}

\begin{frame}{单链表的建立 --- 头插法}
    \textbf{每次将新结点插入到头结点之后}
    \vspace{0.5em}
    \begin{enumerate}
        \item 生成新结点 $s$
        \item $s \to \text{next} = \text{head} \to \text{next}$
        \item $\text{head} \to \text{next} = s$
    \end{enumerate}
    \vspace{0.5em}
    \begin{block}{特点}
        链表顺序与输入顺序相反（逆序建立）
    \end{block}
\end{frame}

\begin{frame}{单链表的建立 --- 尾插法}
    \textbf{每次将新结点插入到链表尾部}
    \vspace{0.5em}
    维护尾指针 $r$ 指向最后一个结点
    \vspace{0.5em}
    \begin{enumerate}
        \item 生成新结点 $s$
        \item $r \to \text{next} = s$
        \item $r = s$
    \end{enumerate}
    \vspace{0.5em}
    \begin{block}{特点}
        链表顺序与输入顺序相同（正序建立）
    \end{block}
\end{frame}

\begin{frame}{单链表的查找}
    \textbf{按位查找}
    \begin{itemize}
        \item 从头结点开始，沿 next 指针移动
        \item 时间复杂度 $O(n)$
    \end{itemize}
    \vspace{0.5em}
    \textbf{按值查找}
    \begin{itemize}
        \item 从第一个元素结点开始依次比较
        \item 时间复杂度 $O(n)$
    \end{itemize}
    \vspace{0.5em}
    \begin{alertblock}{与顺序表对比}
        顺序表按位查找 $O(1)$，链表 $O(n)$
    \end{alertblock}
\end{frame}

\begin{frame}{单链表的插入操作}
    \textbf{后插（已知结点 $p$）}
    \[
    s \to \text{next} = p \to \text{next};\quad p \to \text{next} = s;
    \]
    时间复杂度 $O(1)$
    \vspace{0.5em}
    \textbf{前插（已知结点 $p$）}
    \begin{itemize}
        \item 后插 + 交换数据
        \item 时间复杂度 $O(1)$
    \end{itemize}
    \vspace{0.5em}
    \textbf{按位序插入}
    \begin{itemize}
        \item 先查找第 $i-1$ 个结点
        \item 再后插
        \item 时间复杂度 $O(n)$
    \end{itemize}
\end{frame}

\begin{frame}{单链表的删除操作}
    \textbf{删除 $p$ 的后继 $q$}
    \[
    p \to \text{next} = q \to \text{next};\quad \text{free}(q);
    \]
    时间复杂度 $O(1)$
    \vspace{0.5em}
    \textbf{删除指定结点 $p$}
    \begin{itemize}
        \item 将后继数据复制到 $p$，再删除后继
        \item 时间复杂度 $O(1)$
        \item 注意：$p$ 为尾结点时无效
    \end{itemize}
    \vspace{0.5em}
    \textbf{按位序删除}
    \begin{itemize}
        \item 先查找第 $i-1$ 个结点
        \item 时间复杂度 $O(n)$
    \end{itemize}
\end{frame}

\begin{frame}{例题 --- 单链表就地逆置}
    \textbf{设计算法，对带头结点的单链表就地逆置}
    \[
    O(n), O(1)
    \]
    \vspace{0.3em}
    \textbf{思路}：头插法
    \begin{enumerate}
        \item $\text{head} \to \text{next} = \text{NULL}$
        \item 遍历原链表，依次取下每个结点 $p$
        \item 使用头插法将 $p$ 插入 head 之后
    \end{enumerate}
\end{frame}

\begin{frame}{本节总结}
    \begin{enumerate}
        \item 单链表结点由数据域和指针域组成
        \item 头指针标识链表，头结点简化边界处理
        \item 建立方式：头插法（逆序）、尾插法（正序）
        \item 查找 $O(n)$，已知位置插入删除 $O(1)$
        \item 链表适用于插入删除频繁而查找较少的场景
    \end{enumerate}
\end{frame}

\end{document}


\subsection{1-5 链表进阶}
\documentclass[10pt, aspectratio=169]{beamer}
\usepackage{../beamer_theme_408}

\title{408考研零基础 --- 数据结构}
\subtitle{1-5 链表进阶}
\author{}
\date{}


\begin{document}


\begin{frame}{本节目标}
    \begin{enumerate}
        \item 理解循环链表的结构特征及其与单链表的区别
        \item 理解双向链表的结点结构与插入删除操作
        \item 能够根据实际需求选择顺序表或链表
    \end{enumerate}
\end{frame}

\begin{frame}{循环链表}
    \textbf{特征}：最后一个结点的指针域指向头结点，形成环
    \vspace{0.5em}
    \begin{itemize}
        \item 从任意结点出发可访问所有结点
        \item 判空条件：$\text{head} \to \text{next} == \text{head}$
        \item 遍历终止条件：$p == \text{head}$
    \end{itemize}
    \vspace{0.5em}
    \textbf{尾指针优化}
    \begin{itemize}
        \item 使用尾指针 $\text{rear}$ 标识循环链表
        \item 头结点：$\text{rear} \to \text{next}$
        \item 尾结点：$\text{rear}$ 本身
        \item 可在 $O(1)$ 时间内访问头尾
    \end{itemize}
\end{frame}

\begin{frame}[fragile]{双向链表的结构}
    \textbf{结点包含三个部分}
    \begin{itemize}
        \item 前驱指针 prior
        \item 数据域 data
        \item 后继指针 next
    \end{itemize}
    \vspace{0.5em}
    \begin{lstlisting}[language=C, basicstyle=\ttfamily\small]
typedef struct DNode {
    ElemType data;
    struct DNode *prior, *next;
} DNode, *DLinkList;
    \end{lstlisting}
    \vspace{0.5em}
    \begin{block}{优势}
        支持双向遍历，查找前驱 $O(1)$
    \end{block}
\end{frame}

\begin{frame}{双向链表 --- 插入操作}
    \textbf{在结点 $p$ 之后插入结点 $s$}
    \vspace{0.5em}
    \begin{enumerate}
        \item $s \to \text{next} = p \to \text{next}$
        \item $p \to \text{next} \to \text{prior} = s$（注意判空）
        \item $s \to \text{prior} = p$
        \item $p \to \text{next} = s$
    \end{enumerate}
    \vspace{0.5em}
    \begin{alertblock}{注意}
        若 $p$ 是最后一个结点，$p \to \text{next}$ 为 NULL，步骤 2 跳过
    \end{alertblock}
\end{frame}

\begin{frame}{双向链表 --- 删除操作}
    \textbf{删除 $p$ 的后继结点 $q$}
    \vspace{0.5em}
    \begin{enumerate}
        \item $p \to \text{next} = q \to \text{next}$
        \item $\text{if}(q \to \text{next} \ne \text{NULL}) \; q \to \text{next} \to \text{prior} = p$
        \item $\text{free}(q)$
    \end{enumerate}
    \vspace{0.5em}
    时间复杂度 $O(1)$
\end{frame}

\begin{frame}{顺序表与链表对比}
    \begin{center}
    \begin{tabular}{|c|c|c|}
        \hline
        \textbf{维度} & \textbf{顺序表} & \textbf{链表} \\
        \hline
        存储方式 & 连续 & 离散 \\
        \hline
        随机访问 & $O(1)$ & $O(n)$ \\
        \hline
        插入删除 & $O(n)$ & $O(1)$ \\
        \hline
        存储密度 & 1 & $< 1$ \\
        \hline
        空间分配 & 预分配 & 动态分配 \\
        \hline
    \end{tabular}
    \end{center}
\end{frame}

\begin{frame}{本节总结}
    \begin{enumerate}
        \item 循环链表：尾指针指向头结点，形成环
        \item 双向链表：prior + next 双向指针
        \item 插入删除需修改四个指针，注意边界
        \item 顺序表适合查找密集、规模可预估场景
        \item 链表适合插入删除密集、规模不确定场景
    \end{enumerate}
\end{frame}

\end{document}


\subsection{1-6 栈}
\documentclass[10pt, aspectratio=169]{beamer}

\usetheme{Madrid}
\usecolortheme{dolphin}

\usepackage{ctex}
\usepackage{xeCJK}
\usepackage{amsmath, amssymb}
\usepackage{graphicx}
\usepackage{listings}
\usepackage{tikz}
\usepackage{xcolor}

\title{408考研零基础 --- 数据结构}
\subtitle{1-6 栈}
\author{}
\date{}

\setbeamertemplate{page number in head/foot}{}

\begin{document}


\begin{frame}{本节目标}
    \begin{enumerate}
        \item 理解栈的定义、逻辑特征及其 ADT 定义
        \item 掌握顺序栈和链栈的存储结构与基本操作实现
        \item 理解栈在括号匹配、表达式求值等场景中的应用原理
    \end{enumerate}
\end{frame}

\begin{frame}{栈的定义与 ADT}
    \textbf{栈}：只允许在一端进行插入和删除操作的线性表
    \begin{itemize}
        \item 栈顶 --- 允许操作的一端
        \item 栈底 --- 不允许操作的一端
        \item 原则：后进先出（LIFO）
    \end{itemize}
    \vspace{0.5em}
    \textbf{基本操作}
    \begin{itemize}
        \item InitStack(\&S) --- 初始化
        \item Push(\&S, e) --- 入栈
        \item Pop(\&S, \&e) --- 出栈
        \item GetTop(S, \&e) --- 取栈顶
        \item StackEmpty(S) --- 判栈空
    \end{itemize}
\end{frame}

\begin{frame}[fragile]{顺序栈}
    \begin{lstlisting}[language=C, basicstyle=\ttfamily\small]
#define MaxSize 50
typedef struct {
    ElemType data[MaxSize];
    int top;
} SqStack;
    \end{lstlisting}
    \vspace{0.5em}
    \textbf{核心操作}
    \begin{itemize}
        \item 初始化：$S \to \text{top} = -1$
        \item 入栈：$S \to \text{data}[++S \to \text{top}] = e$
        \item 出栈：$e = S \to \text{data}[S \to \text{top}--]$
        \item 判空：$S \to \text{top} == -1$
    \end{itemize}
    \vspace{0.3em}
    入栈出栈时间复杂度 $O(1)$
\end{frame}

\begin{frame}{共享栈}
    \textbf{两个栈共享同一数组空间}
    \vspace{0.5em}
    \begin{itemize}
        \item 左栈栈底在 0 端，$S_1 \to \text{top} = -1$
        \item 右栈栈底在 MaxSize-1 端，$S_2 \to \text{top} = \text{MaxSize}$
        \item 判满条件：$S_1 \to \text{top} + 1 == S_2 \to \text{top}$
    \end{itemize}
    \vspace{0.5em}
    优点：更有效地利用存储空间
\end{frame}

\begin{frame}{链栈}
    \textbf{操作受限的单链表}
    \vspace{0.5em}
    \begin{itemize}
        \item 入栈：在头结点之后插入
        \item 出栈：删除头结点之后的结点
        \item 无需预分配空间
        \item 不会出现栈满
        \item 时间复杂度 $O(1)$
    \end{itemize}
    \vspace{0.5em}
    \textbf{入栈}
    \[
    s \to \text{next} = S \to \text{next};\quad S \to \text{next} = s;
    \]
    \textbf{出栈}
    \[
    p = S \to \text{next};\quad S \to \text{next} = p \to \text{next};\quad \text{free}(p);
    \]
\end{frame}

\begin{frame}{栈的应用 --- 括号匹配}
    \textbf{算法步骤}
    \begin{enumerate}
        \item 遍历表达式
        \item 遇到左括号 $\to$ 入栈
        \item 遇到右括号 $\to$ 检查栈顶是否为对应左括号
        \begin{itemize}
            \item 匹配 $\to$ 出栈
            \item 不匹配 $\to$ 失败
        \end{itemize}
        \item 遍历结束，栈为空则匹配成功
    \end{enumerate}
\end{frame}

\begin{frame}{栈的应用 --- 表达式求值}
    \textbf{中缀转后缀}
    \begin{itemize}
        \item 操作数 $\to$ 直接输出
        \item 运算符 $\to$ 与栈顶比较优先级
        \begin{itemize}
            \item 优先级高 $\to$ 入栈
            \item 优先级低/相等 $\to$ 出栈并输出
        \end{itemize}
    \end{itemize}
    \vspace{0.5em}
    \textbf{后缀计算}
    \begin{itemize}
        \item 操作数 $\to$ 入栈
        \item 运算符 $\to$ 弹出两个操作数计算，结果入栈
    \end{itemize}
\end{frame}

\begin{frame}{例题 --- 十进制转八进制}
    \textbf{除基取余法}
    \vspace{0.5em}
    十进制 1348 转八进制：
    \begin{itemize}
        \item $1348 \div 8 = 168 \cdots\cdots 4$
        \item $168 \div 8 = 21 \cdots\cdots 0$
        \item $21 \div 8 = 2 \cdots\cdots 5$
        \item $2 \div 8 = 0 \cdots\cdots 2$
    \end{itemize}
    \vspace{0.5em}
    余数入栈 $\to$ 出栈得 2504
    \[
    2 \times 8^3 + 5 \times 8^2 + 0 \times 8^1 + 4 \times 8^0 = 1348
    \]
\end{frame}

\begin{frame}{本节总结}
    \begin{enumerate}
        \item 栈是 LIFO 的受限线性表
        \item 顺序栈：数组实现，$O(1)$，需预分配空间
        \item 链栈：链表实现，$O(1)$，不会栈满
        \item 应用：括号匹配、表达式求值、函数递归
    \end{enumerate}
\end{frame}

\end{document}


\subsection{1-7 队列}
\documentclass[10pt, aspectratio=169]{beamer}

\usetheme{Madrid}
\usecolortheme{dolphin}

\usepackage{ctex}
\usepackage{xeCJK}
\usepackage{amsmath, amssymb}
\usepackage{graphicx}
\usepackage{listings}
\usepackage{tikz}
\usepackage{xcolor}

\title{408考研零基础 --- 数据结构}
\subtitle{1-7 队列}
\author{}
\date{}

\setbeamertemplate{page number in head/foot}{}

\begin{document}


\begin{frame}{本节目标}
    \begin{enumerate}
        \item 理解队列的定义、逻辑特征及其 ADT 定义
        \item 掌握循环队列的存储结构与基本操作
        \item 掌握链队列的存储结构与基本操作
    \end{enumerate}
\end{frame}

\begin{frame}{队列的定义与 ADT}
    \textbf{队列}：只允许在队尾插入、队头删除的线性表
    \begin{itemize}
        \item 队头 --- 删除端
        \item 队尾 --- 插入端
        \item 原则：先进先出（FIFO）
    \end{itemize}
    \vspace{0.5em}
    \textbf{基本操作}
    \begin{itemize}
        \item InitQueue(\&Q) --- 初始化
        \item EnQueue(\&Q, e) --- 入队
        \item DeQueue(\&Q, \&e) --- 出队
        \item GetHead(Q, \&e) --- 取队头
        \item QueueEmpty(Q) --- 判队空
    \end{itemize}
\end{frame}

\begin{frame}{循环队列}
    \textbf{解决假溢出的问题}
    \vspace{0.5em}
    \begin{itemize}
        \item $\text{front}$ --- 队头指针
        \item $\text{rear}$ --- 队尾指针（指向队尾的下一个位置）
    \end{itemize}
    \vspace{0.3em}
    \textbf{模运算实现}
    \[
    \begin{aligned}
    &\text{入队：} Q \to \text{data}[Q \to \text{rear}] = e;\quad Q \to \text{rear} = (Q \to \text{rear} + 1) \% \text{MaxSize}\\
    &\text{出队：} e = Q \to \text{data}[Q \to \text{front}];\quad Q \to \text{front} = (Q \to \text{front} + 1) \% \text{MaxSize}
    \end{aligned}
    \]
    \vspace{0.3em}
    \textbf{队空}：$\text{front} == \text{rear}$ \\
    \textbf{队满}：$(\text{rear} + 1) \% \text{MaxSize} == \text{front}$
\end{frame}

\begin{frame}{循环队列 --- 元素个数}
    \textbf{队列中的元素个数}
    \[
    (\text{rear} - \text{front} + \text{MaxSize}) \% \text{MaxSize}
    \]
    \vspace{0.5em}
    \textbf{区分队空队满的三种方法}
    \begin{enumerate}
        \item 牺牲一个存储单元（常用）
        \item 增加 size 字段
        \item 增加 tag 字段
    \end{enumerate}
\end{frame}

\begin{frame}{链队列}
    \textbf{带队头队尾指针的单链表}
    \vspace{0.5em}
    \begin{lstlisting}[language=C, basicstyle=\ttfamily\small]
typedef struct LinkNode {
    ElemType data;
    struct LinkNode *next;
} LinkNode;
typedef struct {
    LinkNode *front, *rear;
} LinkQueue;
    \end{lstlisting}
    \vspace{0.3em}
    \textbf{特点}
    \begin{itemize}
        \item 不会队满（除非内存耗尽）
        \item 入队出队 $O(1)$
        \item 出队时注意尾指针的更新
    \end{itemize}
\end{frame}

\begin{frame}{队列的应用}
    \textbf{层次遍历}
    \begin{itemize}
        \item 根结点入队 $\to$ 出队访问 $\to$ 子结点入队
    \end{itemize}
    \vspace{0.3em}
    \textbf{缓冲区管理}
    \begin{itemize}
        \item 解决 CPU 与 I/O 设备速度不匹配问题
    \end{itemize}
    \vspace{0.3em}
    \textbf{广度优先搜索 BFS}
    \begin{itemize}
        \item 队列记录待探索的顶点
    \end{itemize}
\end{frame}

\begin{frame}{例题 --- 双栈模拟队列}
    \textbf{利用两个栈模拟队列}
    \vspace{0.5em}
    \begin{itemize}
        \item 栈 A：入队时压入
        \item 栈 B：出队时弹出
    \end{itemize}
    \vspace{0.3em}
    \textbf{入队}：$\text{Push}(A, e)$
    \vspace{0.3em}
    \textbf{出队}
    \begin{itemize}
        \item 若 $B$ 非空，$\text{Pop}(B)$
        \item 若 $B$ 为空，将 $A$ 全部转移到 $B$，再 $\text{Pop}(B)$
    \end{itemize}
    \vspace{0.3em}
    摊还时间复杂度 $O(1)$
\end{frame}

\begin{frame}{本节总结}
    \begin{enumerate}
        \item 队列是 FIFO 的受限线性表
        \item 循环队列用模运算解决假溢出
        \begin{itemize}
            \item 队空 $\text{front} == \text{rear}$
            \item 队满 $(\text{rear}+1)\% \text{MaxSize} == \text{front}$
        \end{itemize}
        \item 链队列使用双指针，入队出队 $O(1)$
        \item 应用：层次遍历、缓冲区、BFS
    \end{enumerate}
\end{frame}

\end{document}


\subsection{1-8 数组与特殊矩阵}
\documentclass[10pt, aspectratio=169]{beamer}
\usepackage{../beamer_theme_408}

\title{408考研零基础 --- 数据结构}
\subtitle{1-8 数组与特殊矩阵}
\author{}
\date{}


\begin{document}


\begin{frame}{本节目标}
    \begin{enumerate}
        \item 理解数组的两种存储方式及地址计算方法
        \item 掌握对称矩阵、三角矩阵、三对角矩阵的压缩存储与地址映射
        \item 理解稀疏矩阵的三元组存储方法
    \end{enumerate}
\end{frame}

\begin{frame}{数组的存储结构}
    \textbf{二维数组 $A[m][n]$}
    \vspace{0.3em}
    \textbf{行优先存储}
    \[
    \text{LOC}(a_{ij}) = \text{LOC}(a_{00}) + (i \times n + j) \times L
    \]
    \vspace{0.3em}
    \textbf{列优先存储}
    \[
    \text{LOC}(a_{ij}) = \text{LOC}(a_{00}) + (j \times m + i) \times L
    \]
    \vspace{0.3em}
    \begin{block}{C 语言}
        采用行优先存储
    \end{block}
\end{frame}

\begin{frame}{对称矩阵的压缩}
    \textbf{$n$ 阶对称矩阵}：$a_{ij} = a_{ji}$
    \vspace{0.3em}
    \textbf{存储下三角 $+$ 主对角线}
    \[
    \text{元素个数} = \frac{n(n+1)}{2}
    \]
    \vspace{0.3em}
    \textbf{地址映射}（$i \ge j$）
    \[
    \text{下标} = \frac{i(i+1)}{2} + j
    \]
    \vspace{0.3em}
    当 $i < j$ 时，访问对称元素 $a_{ji}$
\end{frame}

\begin{frame}{三角矩阵}
    \textbf{下三角矩阵}
    \begin{itemize}
        \item 上三角区域全部为常数 $c$
        \item 存储 $n(n+1)/2 + 1$ 个元素
        \item 地址映射与对称矩阵相同（$i \ge j$）
        \item $i < j$ 时返回常数 $c$
    \end{itemize}
    \vspace{0.3em}
    \textbf{上三角矩阵}
    \begin{itemize}
        \item 按列优先存储上三角区域
    \end{itemize}
\end{frame}

\begin{frame}{三对角矩阵}
    \textbf{非零元素在三条对角线上}
    \[
    |i - j| \le 1
    \]
    \vspace{0.3em}
    \textbf{按行优先存储}
    \[
    \text{下标} = 2i + j
    \]
    \vspace{0.3em}
    \textbf{举例（$5 \times 5$）}
    \[
    a_{00}, a_{01}, a_{10}, a_{11}, a_{12}, a_{21}, a_{22}, a_{23}, a_{32}, a_{33}, a_{34}, a_{43}, a_{44}
    \]
\end{frame}

\begin{frame}{稀疏矩阵 --- 三元组表示}
    \textbf{只存储非零元素}
    \[
    (row, col, value)
    \]
    \vspace{0.3em}
    \begin{center}
    \begin{tabular}{|c|c|c|}
        \hline
        row & col & value \\
        \hline
        0 & 0 & 3 \\
        \hline
        1 & 2 & 5 \\
        \hline
        2 & 0 & 7 \\
        \hline
        \vdots & \vdots & \vdots \\
        \hline
    \end{tabular}
    \end{center}
    \vspace{0.3em}
    \textbf{缺点}：失去随机访问能力
\end{frame}

\begin{frame}{稀疏矩阵 --- 十字链表}
    \textbf{每个非零元素结点}
    \begin{itemize}
        \item 行号、列号、值
        \item 指向同一行下一个非零元素的指针（右指针）
        \item 指向同一列下一个非零元素的指针（下指针）
    \end{itemize}
    \vspace{0.5em}
    优点：支持快速存取任意位置的元素
\end{frame}

\begin{frame}{本节总结}
    \begin{enumerate}
        \item 数组存储：行优先和列优先两种方式
        \item 对称矩阵：存储 $n(n+1)/2$ 个元素
        \item 三角矩阵：下三角/上三角 $+$ 常数
        \item 三对角矩阵：存储三条对角线
        \item 稀疏矩阵：三元组或十字链表
    \end{enumerate}
\end{frame}

\end{document}


\subsection{1-9 树的基本概念}
\documentclass[10pt, aspectratio=169]{beamer}

\usetheme{Madrid}
\usecolortheme{dolphin}

\usepackage{ctex}
\usepackage{xeCJK}
\usepackage{amsmath, amssymb}
\usepackage{graphicx}
\usepackage{listings}
\usepackage{tikz}
\usepackage{xcolor}

\title{408考研零基础 --- 数据结构}
\subtitle{1-9 树的基本概念}
\author{}
\date{}

\setbeamertemplate{page number in head/foot}{}

\begin{document}


\begin{frame}{本节目标}
    \begin{enumerate}
        \item 理解树的递归定义与基本术语
        \item 掌握树的性质（结点数、度数、深度之间的关系）
        \item 理解树的四种存储结构及特点
    \end{enumerate}
\end{frame}

\begin{frame}{树的递归定义}
    \textbf{树}：$n$ 个结点的有限集合（$n \ge 0$）
    \vspace{0.3em}
    \begin{itemize}
        \item $n = 0$ 时为空树
        \item 非空树有且仅有一个根结点
        \item 其余结点分为 $m$ 个互不相交的有限集合，每个集合本身是一棵树（子树）
    \end{itemize}
    \vspace{0.5em}
    \begin{block}{特征}
        递归定义，树由根结点和若干子树构成，每棵子树又是一棵树
    \end{block}
\end{frame}

\begin{frame}{树的术语}
    \begin{itemize}
        \item \textbf{度}：结点拥有的子树个数
        \item \textbf{树的度}：所有结点度的最大值
        \item \textbf{叶结点}：度为 0 的结点
        \item \textbf{分支结点}：度 $> 0$ 的结点
        \item \textbf{深度/高度}：结点的最大层次数
        \item \textbf{双亲} $\leftrightarrow$ \textbf{孩子}：上下层关系
        \item \textbf{兄弟}：同一双亲的孩子
        \item \textbf{祖先} $\leftrightarrow$ \textbf{子孙}：路径上的前后关系
    \end{itemize}
\end{frame}

\begin{frame}{树的性质}
    \textbf{性质一}
    \[
    \text{结点数} = \text{所有结点的度数之和} + 1
    \]
    \vspace{0.3em}
    \textbf{性质二}：度为 $m$ 的树第 $i$ 层最多有
    \[
    m^{i-1} \text{ 个结点}
    \]
    \vspace{0.3em}
    \textbf{性质三}：高度为 $h$ 的 $m$ 叉树最多有
    \[
    \frac{m^h - 1}{m - 1} \text{ 个结点}
    \]
    \vspace{0.3em}
    \textbf{性质四}：$n$ 个结点的 $m$ 叉树最小高度
    \[
    \lceil \log_m(n(m-1) + 1) \rceil
    \]
\end{frame}

\begin{frame}{树的存储结构}
    \textbf{双亲表示法}
    \begin{itemize}
        \item 数组存储结点，每个结点含数据域和双亲指针
        \item 找双亲 $O(1)$，找孩子需遍历
    \end{itemize}
    \vspace{0.3em}
    \textbf{孩子表示法}
    \begin{itemize}
        \item 每个结点维护指向孩子链表的指针
        \item 找孩子方便，找双亲需遍历
    \end{itemize}
    \vspace{0.3em}
    \textbf{孩子兄弟表示法}
    \begin{itemize}
        \item 每个结点含 firstchild 和 nextsibling 指针
        \item 可将树转换为二叉树
    \end{itemize}
\end{frame}

\begin{frame}{例题 --- 求叶结点数}
    \textbf{一棵度为 4 的树中，度 1、2、3、4 的结点数分别为 4、2、1、1，求叶结点数 $n_0$}
    \vspace{0.5em}
    \textbf{解答}
    \[
    \begin{aligned}
    n &= n_0 + 4 + 2 + 1 + 1 = n_0 + 8\\
    \text{总度数} &= 0 \times n_0 + 1 \times 4 + 2 \times 2 + 3 \times 1 + 4 \times 1 = 15\\
    n &= 15 + 1 = 16\\
    n_0 + 8 &= 16 \implies n_0 = 8
    \end{aligned}
    \]
    \vspace{0.3em}
    叶结点数为 \textbf{8}
\end{frame}

\begin{frame}{本节总结}
    \begin{enumerate}
        \item 树是递归定义的分层结构
        \item 术语：度、叶结点、深度、双亲、孩子、兄弟
        \item 结点数 $=$ 度数之和 $+1$
        \item 存储结构：双亲、孩子、孩子兄弟，各有适用场景
    \end{enumerate}
\end{frame}

\end{document}


\subsection{1-10 二叉树}
\documentclass[10pt, aspectratio=169]{beamer}

\usetheme{Madrid}
\usecolortheme{dolphin}

\usepackage{ctex}
\usepackage{xeCJK}
\usepackage{amsmath, amssymb}
\usepackage{graphicx}
\usepackage{listings}
\usepackage{tikz}
\usepackage{xcolor}

\title{408考研零基础 --- 数据结构}
\subtitle{1-10 二叉树}
\author{}
\date{}

\setbeamertemplate{page number in head/foot}{}

\begin{document}


\begin{frame}{本节目标}
    \begin{enumerate}
        \item 理解二叉树的递归定义及其与度为 2 的树的区别
        \item 掌握二叉树的重要性质（$n_0 = n_2 + 1$）
        \item 掌握满二叉树、完全二叉树、二叉排序树、平衡二叉树的概念与特征
    \end{enumerate}
\end{frame}

\begin{frame}{二叉树的递归定义}
    \textbf{二叉树}：$n$ 个结点的有限集合
    \begin{itemize}
        \item $n = 0$ 时为空二叉树
        \item 非空二叉树有且仅有一个根结点
        \item 其余结点分为两个互不相交的集合：左子树和右子树
    \end{itemize}
    \vspace{0.5em}
    \begin{alertblock}{与度为 2 的树的区别}
        \begin{itemize}
            \item 二叉树明确区分左右子树
            \item 度为 2 的树不区分左右
            \item 二叉树允许只有一个子结点时指明左右
        \end{itemize}
    \end{alertblock}
\end{frame}

\begin{frame}{二叉树的性质}
    \textbf{性质一}：第 $i$ 层最多 $2^{i-1}$ 个结点
    \vspace{0.3em}
    \textbf{性质二}：深度为 $h$ 的二叉树最多 $2^h - 1$ 个结点
    \vspace{0.3em}
    \textbf{性质三（核心）}
    \[
    n_0 = n_2 + 1
    \]
    叶结点数 $=$ 度为 2 的结点数 $+ 1$
    \vspace{0.3em}
    \textbf{证明}
    \[
    \begin{aligned}
    \text{总度数} &= n - 1 = n_1 + 2n_2\\
    n &= n_0 + n_1 + n_2\\
    \therefore\; n_0 &= n_2 + 1
    \end{aligned}
    \]
\end{frame}

\begin{frame}{满二叉树}
    \textbf{定义}：深度为 $h$，结点数为 $2^h - 1$
    \vspace{0.5em}
    \textbf{特征}
    \begin{itemize}
        \item 所有叶结点在第 $h$ 层
        \item 每个分支结点都有两棵非空子树
        \item 结点数达到最大值
    \end{itemize}
\end{frame}

\begin{frame}{完全二叉树}
    \textbf{定义}：结点编号与满二叉树 1 至 $n$ 对应
    \vspace{0.5em}
    \textbf{特征}
    \begin{itemize}
        \item 叶结点只可能在最后两层
        \item 最多一个度为 1 的结点（且只有左子树）
        \item 结点 $i$ 的左孩子为 $2i$，右孩子为 $2i+1$
        \item 结点 $i$ 的双亲为 $\lfloor i/2 \rfloor$
    \end{itemize}
    \vspace{0.5em}
    \begin{block}{注意}
        完全二叉树非常适合数组存储
    \end{block}
\end{frame}

\begin{frame}{二叉排序树与平衡二叉树}
    \textbf{二叉排序树（BST）}
    \begin{itemize}
        \item 左子树所有结点 $<$ 根 $<$ 右子树所有结点
        \item 左右子树本身也是二叉排序树
    \end{itemize}
    \vspace{0.3em}
    \textbf{平衡二叉树（AVL）}
    \begin{itemize}
        \item 任意结点左右子树高度差 $\le 1$
        \item 查找效率稳定在 $O(\log n)$
    \end{itemize}
\end{frame}

\begin{frame}{二叉树的存储结构}
    \textbf{顺序存储}
    \begin{itemize}
        \item 按完全二叉树编号存储
        \item 适用于完全二叉树和满二叉树
        \item 普通二叉树可能浪费空间
    \end{itemize}
    \vspace{0.3em}
    \textbf{链式存储（二叉链表）}
    \begin{lstlisting}[language=C, basicstyle=\ttfamily\small]
typedef struct BiTNode {
    ElemType data;
    struct BiTNode *lchild, *rchild;
} BiTNode, *BiTree;
    \end{lstlisting}
    \vspace{0.3em}
    $n$ 个结点有 $n+1$ 个空指针域
\end{frame}

\begin{frame}{例题 --- 完全二叉树结点数}
    \textbf{完全二叉树有 1000 个结点，求 $n_0, n_1, n_2$}
    \vspace{0.5em}
    \textbf{解答}
    \[
    \begin{aligned}
    n_0 &= n_2 + 1\\
    n_0 + n_1 + n_2 &= 1000 \implies 2n_2 + n_1 = 999
    \end{aligned}
    \]
    $n_1$ 只能为 0 或 1
    \begin{itemize}
        \item $n_1 = 0 \implies 2n_2 = 999$，不成立
        \item $n_1 = 1 \implies 2n_2 = 998 \implies n_2 = 499$
    \end{itemize}
    \[
    n_0 = 500,\; n_1 = 1,\; n_2 = 499
    \]
\end{frame}

\begin{frame}{本节总结}
    \begin{enumerate}
        \item 二叉树每个结点至多两棵子树，区分左右
        \item 核心性质：$n_0 = n_2 + 1$
        \item 满二叉树：$2^h - 1$ 个结点
        \item 完全二叉树：适合数组存储
        \item BST 和 AVL 是两种重要的特殊二叉树
        \item 链式存储：$n$ 个结点 $n+1$ 个空指针
    \end{enumerate}
\end{frame}

\end{document}


\subsection{1-11 二叉树的遍历}
\documentclass[10pt, aspectratio=169]{beamer}
\usepackage{../beamer_theme_408}

\title{408考研零基础 --- 数据结构}
\subtitle{1-11 二叉树的遍历}
\author{}
\date{}


\begin{document}


\begin{frame}{本节目标}
    \begin{enumerate}
        \item 理解先序、中序、后序遍历的递归定义与访问顺序
        \item 掌握非递归遍历算法的实现
        \item 掌握层序遍历的队列实现，能根据遍历序列还原二叉树
    \end{enumerate}
\end{frame}

\begin{frame}{递归遍历}
    \textbf{先序遍历（根左右）}
    \[
    \text{根} \to \text{左子树} \to \text{右子树}
    \]
    \textbf{中序遍历（左根右）}
    \[
    \text{左子树} \to \text{根} \to \text{右子树}
    \]
    \textbf{后序遍历（左右根）}
    \[
    \text{左子树} \to \text{右子树} \to \text{根}
    \]
    \vspace{0.5em}
    \begin{block}{核心区别}
        根结点的访问时机决定遍历方式名称
    \end{block}
\end{frame}

\begin{frame}[fragile]{递归算法实现}
    \textbf{先序遍历}
    \begin{lstlisting}[language=C, basicstyle=\ttfamily\small]
void PreOrder(BiTree T) {
    if(T != NULL) {
        visit(T);
        PreOrder(T->lchild);
        PreOrder(T->rchild);
    }
}
    \end{lstlisting}
    \vspace{0.3em}
    \textbf{中序和后序仅调整 visit 位置}
    \vspace{0.3em}
    时间复杂度 $O(n)$，空间复杂度 $O(h)$
\end{frame}

\begin{frame}{非递归遍历 --- 先序}
    \textbf{使用栈模拟}
    \vspace{0.3em}
    \begin{enumerate}
        \item 根结点入栈
        \item 出栈 $p$，访问 $p$
        \item 右孩子入栈（先右后左，保证左先出）
        \item 左孩子入栈
        \item 重复直到栈空
    \end{enumerate}
\end{frame}

\begin{frame}{非递归遍历 --- 中序}
    \textbf{使用栈模拟}
    \vspace{0.3em}
    \begin{enumerate}
        \item $p$ 指向根
        \item 若 $p$ 非空，$p$ 入栈，$p = p \to \text{lchild}$
        \item 若 $p$ 为空且栈非空，出栈 $p$，访问 $p$
        \item $p = p \to \text{rchild}$
        \item 重复直到栈空且 $p$ 为空
    \end{enumerate}
\end{frame}

\begin{frame}{层序遍历}
    \textbf{使用队列}
    \vspace{0.3em}
    \begin{enumerate}
        \item 根结点入队
        \item 出队 $p$，访问 $p$
        \item 左孩子入队
        \item 右孩子入队
        \item 重复直到队空
    \end{enumerate}
    \vspace{0.5em}
    \textbf{时间复杂度} $O(n)$，\textbf{空间复杂度} $O(n)$
    \vspace{0.3em}
    \textbf{应用}：求最大宽度、判断完全二叉树
\end{frame}

\begin{frame}{由遍历序列还原二叉树}
    \textbf{先序 + 中序 $\implies$ 唯一二叉树}
    \vspace{0.3em}
    \textbf{原理}
    \begin{enumerate}
        \item 先序第一个为根
        \item 中序中根左边为左子树，右边为右子树
        \item 递归构建
    \end{enumerate}
    \vspace{0.3em}
    \textbf{例题}
    \begin{itemize}
        \item 先序：A B D E C F
        \item 中序：D B E A F C
    \end{itemize}
    根 A，左子树 (B D E) / (D B E)，右子树 (C F) / (F C)
\end{frame}

\begin{frame}{本节总结}
    \begin{enumerate}
        \item 先序（根左右）、中序（左根右）、后序（左右根）
        \item 递归遍历 $O(n)$ 时间，$O(h)$ 空间
        \item 非递归遍历用栈模拟递归
        \item 层序遍历用队列实现
        \item 先序/后序 + 中序可唯一还原二叉树
    \end{enumerate}
\end{frame}

\end{document}


\subsection{1-12 线索二叉树}
\documentclass[10pt, aspectratio=169]{beamer}

\usetheme{Madrid}
\usecolortheme{dolphin}

\usepackage{ctex}
\usepackage{xeCJK}
\usepackage{amsmath, amssymb}
\usepackage{graphicx}
\usepackage{listings}
\usepackage{tikz}
\usepackage{xcolor}

\title{408考研零基础 --- 数据结构}
\subtitle{1-12 线索二叉树}
\author{}
\date{}

\setbeamertemplate{page number in head/foot}{}

\begin{document}


\begin{frame}{本节目标}
    \begin{enumerate}
        \item 理解线索化的动机与原理
        \item 掌握二叉树的中序线索化算法
        \item 掌握线索二叉树中查找前驱与后继的方法
    \end{enumerate}
\end{frame}

\begin{frame}{线索化的动机}
    \textbf{$n$ 个结点的二叉链表}
    \begin{itemize}
        \item 共 $2n$ 个指针域
        \item $n-1$ 个指向孩子
        \item \textbf{$n+1$ 个空指针域}
    \end{itemize}
    \vspace{0.5em}
    \textbf{线索化思想}
    \begin{itemize}
        \item 空左指针 $\to$ 指向前驱
        \item 空右指针 $\to$ 指向后继
        \item 增加标志位区分孩子/线索
    \end{itemize}
\end{frame}

\begin{frame}{线索二叉树的存储结构}
    \begin{lstlisting}[language=C, basicstyle=\ttfamily\small]
typedef struct ThreadNode {
    ElemType data;
    struct ThreadNode *lchild, *rchild;
    int ltag, rtag;
} ThreadNode, *ThreadTree;
    \end{lstlisting}
    \vspace{0.5em}
    \textbf{标志位含义}
    \begin{itemize}
        \item $\text{ltag} = 0$：lchild 指向左孩子
        \item $\text{ltag} = 1$：lchild 指向前驱
        \item $\text{rtag} = 0$：rchild 指向右孩子
        \item $\text{rtag} = 1$：rchild 指向后继
    \end{itemize}
\end{frame}

\begin{frame}{中序线索化算法}
    \textbf{基于中序遍历，使用 pre 指针}
    \vspace{0.3em}
    \begin{enumerate}
        \item 线索化左子树
        \item 处理当前结点 $p$
        \begin{itemize}
            \item 若 $p \to \text{lchild} == \text{NULL}$，指向 pre
            \item 若 $\text{pre} \to \text{rchild} == \text{NULL}$，指向 $p$
        \end{itemize}
        \item $\text{pre} = p$
        \item 线索化右子树
    \end{enumerate}
    \vspace{0.3em}
    时间复杂度 $O(n)$，空间复杂度 $O(1)$
\end{frame}

\begin{frame}{中序线索二叉树的遍历}
    \textbf{找中序后继}
    \begin{itemize}
        \item $\text{rtag} = 1$：直接 $p \to \text{rchild}$
        \item $\text{rtag} = 0$：右子树中最左边的结点
    \end{itemize}
    \vspace{0.3em}
    \textbf{找中序前驱}
    \begin{itemize}
        \item $\text{ltag} = 1$：直接 $p \to \text{lchild}$
        \item $\text{ltag} = 0$：左子树中最右边的结点
    \end{itemize}
    \vspace{0.3em}
    \textbf{遍历算法}
    \begin{enumerate}
        \item 找到中序第一个结点（沿左走到 ltag=1）
        \item 循环调用找后继，访问每个结点
    \end{enumerate}
    时间复杂度 $O(n)$，空间复杂度 $O(1)$
\end{frame}

\begin{frame}{本节总结}
    \begin{enumerate}
        \item 线索化利用 $n+1$ 个空指针记录前驱后继
        \item 中序线索化基于中序遍历，使用 pre 指针
        \item 线索二叉树遍历无需递归或栈
        \item 时间 $O(n)$，空间 $O(1)$
    \end{enumerate}
\end{frame}

\end{document}


\subsection{1-13 树、森林与二叉树}
\documentclass[10pt, aspectratio=169]{beamer}
\usepackage{../beamer_theme_408}

\title{408考研零基础 --- 数据结构}
\subtitle{1-13 树、森林与二叉树}
\author{}
\date{}


\begin{document}


\begin{frame}{本节目标}
    \begin{enumerate}
        \item 掌握树转换为二叉树的方法
        \item 掌握森林转换为二叉树的方法及还原
        \item 理解树和森林的遍历方式及其与二叉树遍历的对应关系
    \end{enumerate}
\end{frame}

\begin{frame}{树转换为二叉树}
    \textbf{孩子兄弟表示法}
    \vspace{0.3em}
    \textbf{转换步骤}
    \begin{enumerate}
        \item 兄弟之间加连线
        \item 只保留与第一个孩子的连线，删除与其他孩子的连线
        \item 以根为轴顺时针旋转 $45^\circ$
    \end{enumerate}
    \vspace{0.3em}
    \textbf{转换结果}
    \begin{itemize}
        \item 左指针 $\to$ 第一个孩子
        \item 右指针 $\to$ 下一个兄弟
        \item 根结点无右子树
    \end{itemize}
\end{frame}

\begin{frame}{森林转换为二叉树}
    \textbf{转换步骤}
    \begin{enumerate}
        \item 每棵树分别转换为二叉树
        \item 用右指针连接各棵树的根
    \end{enumerate}
    \vspace{0.5em}
    \textbf{二叉树 $\to$ 森林}
    \begin{itemize}
        \item 根为第一棵树根
        \item 左子树为第一棵树的子结点结构
        \item 右指针指向第二棵树根
    \end{itemize}
\end{frame}

\begin{frame}{树和森林的遍历}
    \textbf{树的先根遍历} $\leftrightarrow$ \textbf{二叉树的先序遍历}
    \vspace{0.3em}
    \textbf{树的后根遍历} $\leftrightarrow$ \textbf{二叉树的中序遍历}
    \vspace{0.5em}
    \textbf{森林的先序遍历}
    \begin{itemize}
        \item 访问第一棵树根
        \item 先序遍历第一棵树的子树森林
        \item 先序遍历剩余森林
    \end{itemize}
    \vspace{0.3em}
    \textbf{森林的中序遍历}
    \begin{itemize}
        \item 中序遍历第一棵树的子树森林
        \item 访问第一棵树根
        \item 中序遍历剩余森林
    \end{itemize}
\end{frame}

\begin{frame}{例题}
    \textbf{森林 F 对应二叉树 B 有 $m$ 个结点，B 的根为 $p$，$p$ 的右孩子结点数为 $n$，求森林中树的棵数}
    \vspace{0.5em}
    \textbf{分析}
    \begin{itemize}
        \item 根 $p$ 对应第一棵树
        \item $p$ 的右孩子对应第二棵树根
        \item 右孩子的右孩子对应第三棵树根
        \item 以此类推
    \end{itemize}
    \vspace{0.5em}
    \[
    \boxed{\text{树的棵数} = n + 1}
    \]
\end{frame}

\begin{frame}{本节总结}
    \begin{enumerate}
        \item 树 $\to$ 二叉树：孩子兄弟表示法
        \item 森林 $\to$ 二叉树：每棵树转换后根用右指针连接
        \item 树的先根遍历 $\leftrightarrow$ 二叉树先序遍历
        \item 树的后根遍历 $\leftrightarrow$ 二叉树中序遍历
        \item 孩子兄弟表示法统一了树、森林、二叉树的存储
    \end{enumerate}
\end{frame}

\end{document}


\subsection{1-14 哈夫曼树与并查集}
\documentclass[10pt, aspectratio=169]{beamer}
\usepackage{../beamer_theme_408}

\title{408考研零基础 --- 数据结构}
\subtitle{1-14 哈夫曼树与并查集}
\author{}
\date{}


\begin{document}


\begin{frame}{本节目标}
    \begin{enumerate}
        \item 理解带权路径长度 WPL 的定义，掌握哈夫曼树的构造算法
        \item 理解哈夫曼编码的原理
        \item 理解并查集的基本概念和核心操作
    \end{enumerate}
\end{frame}

\begin{frame}{哈夫曼树的基本概念}
    \textbf{路径长度}：路径上的分支数目
    \vspace{0.3em}
    \textbf{结点的带权路径长度}
    \[
    \text{路径长度} \times \text{结点的权}
    \]
    \vspace{0.3em}
    \textbf{树的带权路径长度（WPL）}
    \[
    \text{WPL} = \sum_{i=1}^{n} w_i \times l_i
    \]
    \vspace{0.3em}
    \textbf{哈夫曼树}：WPL 最小的二叉树（最优二叉树）
    \vspace{0.3em}
    \begin{block}{特点}
        权值越大的叶结点离根越近
    \end{block}
\end{frame}

\begin{frame}{哈夫曼树构造算法}
    \textbf{给定 $n$ 个权值 $\{w_1, w_2, \ldots, w_n\}$}
    \vspace{0.3em}
    \begin{enumerate}
        \item 每个权值视为一棵只有根结点的二叉树
        \item 选择两棵权值最小的树合并
        \item 新根权值 $=$ 两子树根权值之和
        \item 删除这两棵树，加入新树
        \item 重复步骤 2-4 直到只剩一棵树
    \end{enumerate}
    \vspace{0.3em}
    $n$ 个叶结点的哈夫曼树共有 $2n - 1$ 个结点
\end{frame}

\begin{frame}{哈夫曼编码}
    \textbf{可变长度前缀编码}
    \vspace{0.3em}
    \textbf{生成方法}
    \begin{itemize}
        \item 字符频率作为叶结点权值
        \item 构造哈夫曼树
        \item 左分支编码 0，右分支编码 1
        \item 路径上的 0/1 序列即为编码
    \end{itemize}
    \vspace{0.3em}
    \textbf{特征}
    \begin{itemize}
        \item 前缀编码：无歧义
        \item 频率越高的字符编码越短
        \item 平均编码长度最小
    \end{itemize}
    \vspace{0.3em}
    \textbf{应用}：ZIP 压缩、JPEG 图像压缩
\end{frame}

\begin{frame}{并查集}
    \textbf{处理不相交集合的合并与查询}
    \vspace{0.3em}
    \textbf{两个核心操作}
    \begin{itemize}
        \item \textbf{Find}：确定元素所属集合（返回根结点）
        \item \textbf{Union}：合并两个集合
    \end{itemize}
    \vspace{0.3em}
    \textbf{存储结构}：双亲表示法的树
    \[
    \text{parent}[i] = 
    \begin{cases}
    i, & i \text{ 为根}\\
    \text{parent}, & \text{其他}
    \end{cases}
    \]
\end{frame}

\begin{frame}{并查集的优化}
    \textbf{按秩合并}
    \begin{itemize}
        \item 高度小的树合并到高度大的树
        \item 避免树高度增长
    \end{itemize}
    \vspace{0.3em}
    \textbf{路径压缩}
    \begin{itemize}
        \item Find 操作中将路径上所有结点直接指向根
        \item 后续查找 $O(1)$
    \end{itemize}
    \vspace{0.3em}
    优化后时间复杂度 $O(\alpha(n))$，$\alpha(n)$ 为阿克曼反函数
\end{frame}

\begin{frame}{本节总结}
    \begin{enumerate}
        \item 哈夫曼树：WPL 最小的二叉树
        \item 构造：每次合并权值最小的两棵树
        \item 哈夫曼编码：最优前缀编码，用于数据压缩
        \item 并查集：Find + Union 两个操作
        \item 优化：路径压缩 + 按秩合并 $\to O(\alpha(n))$
    \end{enumerate}
\end{frame}

\end{document}


\subsection{1-15 图的基本概念与存储}
\documentclass[10pt, aspectratio=169]{beamer}
\usetheme{Madrid}
\usecolortheme{dolphin}
\usepackage{ctex}
\usepackage{xeCJK}
\usepackage{amsmath,amssymb}
\usepackage{listings}
\usepackage{tikz}
\usepackage{xcolor}
\title{408考研零基础 --- 数据结构}
\subtitle{1-15 图的基本概念与存储}
\author{}
\date{}
\setbeamertemplate{page number in head/foot}{}
\begin{document}
\frame{\titlepage}
\begin{frame}{本节目标}
\begin{enumerate}
\item 掌握图的定义和术语
\item 理解邻接矩阵和邻接表的存储结构
\item 能够根据实际需求选择合适的存储结构
\end{enumerate}
\end{frame}
\begin{frame}{图的定义与术语}
\textbf{图} $G=(V,E)$
\begin{itemize}
\item $V$ --- 顶点的非空有限集合
\item $E$ --- 边的有限集合
\end{itemize}
\vspace{0.3em}
\begin{tabular}{|c|c|}
\hline 术语 & 定义 \\ \hline
有向图 & 边有方向（弧） \\ \hline
无向图 & 边无方向 \\ \hline
度 & 与顶点相关联的边数 \\ \hline
入度/出度 & 有向图中指向/指出的边数 \\ \hline
连通分量 & 无向图中的极大连通子图 \\ \hline
强连通分量 & 有向图中的极大强连通子图 \\ \hline
\end{tabular}
\end{frame}
\begin{frame}{完全图}
\textbf{无向完全图} $|E|=n(n-1)/2$
\vspace{0.3em}
\textbf{有向完全图} $|E|=n(n-1)$
\vspace{0.5em}
\textbf{稀疏图}：边数很少\quad \textbf{稠密图}：边数接近完全图
\end{frame}
\begin{frame}{邻接矩阵}
$n$ 个顶点的 $n\times n$ 方阵，$A[i][j]=1$ 表示存在边
\vspace{0.3em}
\textbf{特点}
\begin{itemize}
\item 无向图对称，有向图不一定对称
\item 判边 $O(1)$，空间 $O(n^2)$
\item 适合稠密图
\end{itemize}
顶点 $i$ 的度：无向图 $\to$ 第 $i$ 行之和
\end{frame}
\begin{frame}{邻接表}
顶点表（数组）+ 边表（链表）
\vspace{0.3em}
\textbf{特点}
\begin{itemize}
\item 空间 $O(n+e)$，适合稀疏图
\item 判边 $O(\text{degree}(v))$
\item 遍历邻接顶点更高效
\end{itemize}
\vspace{0.3em}
\begin{tabular}{|c|c|c|}
\hline 对比 & 邻接矩阵 & 邻接表 \\ \hline
空间 & $O(n^2)$ & $O(n+e)$ \\ \hline
判边 & $O(1)$ & $O(\text{degree})$ \\ \hline
\end{tabular}
\end{frame}
\begin{frame}{本节总结}
\begin{enumerate}
\item 图分为有向图和无向图
\item 邻接矩阵空间 $O(n^2)$，适合稠密图，判边 $O(1)$
\item 邻接表空间 $O(n+e)$，适合稀疏图
\end{enumerate}
\vspace{1em}
\begin{center}
\textcolor{blue}{\textbf{下节预告：图的遍历}}
\end{center}
\end{frame}
\end{document}


\subsection{1-16 图的遍历}
\documentclass[10pt, aspectratio=169]{beamer}
\usetheme{Madrid}
\usecolortheme{dolphin}
\usepackage{ctex}
\usepackage{xeCJK}
\usepackage{amsmath,amssymb}
\usepackage{listings}
\usepackage{tikz}
\usepackage{xcolor}
\title{408考研零基础 --- 数据结构}
\subtitle{1-16 图的遍历}
\author{}
\date{}
\setbeamertemplate{page number in head/foot}{}
\begin{document}
\frame{\titlepage}
\begin{frame}{本节目标}
\begin{enumerate}
\item 理解 DFS 的原理与实现
\item 理解 BFS 的原理与实现
\item 掌握 DFS 和 BFS 在连通性、生成树中的应用
\end{enumerate}
\end{frame}
\begin{frame}{深度优先搜索 DFS}
\textbf{策略}：沿路径深入探索，回溯继续
\vspace{0.3em}
\textbf{递归算法}
\begin{itemize}
\item 访问 $v$，标记 visited
\item 取 $v$ 的未访问邻接顶点 $w$，递归 DFS($w$)
\item 所有邻接顶点已访问则回溯
\end{itemize}
\vspace{0.3em}
\textbf{时间复杂度}
\begin{itemize}
\item 邻接表：$O(n+e)$
\item 邻接矩阵：$O(n^2)$
\end{itemize}
空间复杂度 $O(n)$（递归栈）
\end{frame}
\begin{frame}{广度优先搜索 BFS}
\textbf{策略}：按距离递增逐层访问
\vspace{0.3em}
\textbf{队列实现}
\begin{itemize}
\item 访问 $v$，入队
\item 出队 $u$，访问 $u$ 的所有未访问邻接顶点并入队
\item 重复直到队空
\end{itemize}
\vspace{0.3em}
\textbf{时间复杂度}
\begin{itemize}
\item 邻接表：$O(n+e)$
\item 邻接矩阵：$O(n^2)$
\end{itemize}
空间复杂度 $O(n)$（队列）
\vspace{0.3em}
\textbf{应用}：无权图单源最短路径
\end{frame}
\begin{frame}{图的连通性判定}
一次 DFS/BFS 只能访问一个连通分量
\vspace{0.3em}
\textbf{遍历所有顶点}
\begin{itemize}
\item 对每个未访问顶点启动一次遍历
\item 遍历次数 $=$ 连通分量数
\end{itemize}
\vspace{0.5em}
\textbf{生成树}
\begin{itemize}
\item DFS 路径 $\to$ 深度优先生成树
\item BFS 路径 $\to$ 广度优先生成树
\item 包含全部 $n$ 个顶点，$n-1$ 条边
\end{itemize}
\end{frame}
\begin{frame}{本节总结}
\begin{enumerate}
\item DFS：递归/栈，沿路径深入，$O(n+e)$
\item BFS：队列，按层访问，$O(n+e)$
\item 一次遍历一个连通分量
\item 遍历路径构成生成树
\end{enumerate}
\vspace{1em}
\begin{center}
\textcolor{blue}{\textbf{下节预告：查找基础}}
\end{center}
\end{frame}
\end{document}


\subsection{1-17 查找基础}
\documentclass[10pt, aspectratio=169]{beamer}
\usetheme{Madrid}
\usecolortheme{dolphin}
\usepackage{ctex}
\usepackage{xeCJK}
\usepackage{amsmath,amssymb}
\usepackage{listings}
\usepackage{tikz}
\usepackage{xcolor}
\title{408考研零基础 --- 数据结构}
\subtitle{1-17 查找基础}
\author{}
\date{}
\setbeamertemplate{page number in head/foot}{}
\begin{document}
\frame{\titlepage}
\begin{frame}{本节目标}
\begin{enumerate}
\item 掌握顺序查找的算法与性能
\item 掌握折半查找的算法与判定树
\item 理解分块查找和 B 树基本概念
\end{enumerate}
\end{frame}
\begin{frame}{查找的基本概念}
\textbf{平均查找长度 ASL}
\[
\text{ASL} = \sum_{i=1}^{n} p_i \times c_i
\]
\vspace{0.3em}
\textbf{分类}
\begin{itemize}
\item 静态查找表：只查找，不修改
\item 动态查找表：查找同时插入/删除
\end{itemize}
\end{frame}
\begin{frame}{顺序查找}
\textbf{算法}：从表一端开始依次比较
\vspace{0.3em}
\textbf{性能}
\begin{itemize}
\item 查找成功 ASL：$(n+1)/2$
\item 查找失败比较次数：$n+1$
\item 时间复杂度 $O(n)$
\end{itemize}
\vspace{0.3em}
\textbf{优化}：设置哨兵减少边界判断
\vspace{0.3em}
适用于无序表或小规模数据
\end{frame}
\begin{frame}{折半查找}
\textbf{适用条件}：有序的顺序表
\vspace{0.3em}
\textbf{步骤}
\begin{itemize}
\item low, high 指向两端，mid $=$ (low+high)/2
\item 关键字 $=$ mid 元素 $\to$ 成功
\item 关键字 $<$ mid $\to$ high $=$ mid$-1$
\item 关键字 $>$ mid $\to$ low $=$ mid$+1$
\item low $>$ high $\to$ 失败
\end{itemize}
\vspace{0.3em}
\textbf{性能}：ASL $\approx \log_2(n+1)-1$，$O(\log n)$
\end{frame}
\begin{frame}{分块查找}
\textbf{块间有序，块内无序}
\vspace{0.3em}
\textbf{步骤}
\begin{enumerate}
\item 索引表中折半查找确定所属块
\item 块内顺序查找
\end{enumerate}
\vspace{0.3em}
\textbf{性能}：$s=\sqrt{n}$ 时 ASL 最小 $\approx \sqrt{n}+1$
\vspace{0.3em}
适合动态插入场景
\end{frame}
\begin{frame}{B 树的基本概念}
$m$ 阶 B 树
\begin{itemize}
\item 每个结点至多 $m$ 棵子树（$m-1$ 个关键字）
\item 根至少 2 棵子树
\item 非叶结点至少 $\lceil m/2 \rceil$ 棵子树
\item 所有叶结点在同一层
\end{itemize}
\vspace{0.3em}
\textbf{查找}：$O(\log_m n)$，降低树高减少磁盘 I/O
\vspace{0.3em}
广泛应用于文件系统和数据库索引
\end{frame}
\begin{frame}{本节总结}
\begin{enumerate}
\item 顺序查找：$O(n)$，ASL $(n+1)/2$
\item 折半查找：$O(\log n)$，要求有序顺序表
\item 分块查找：索引 + 顺序，适合动态场景
\item B 树：多路平衡查找，$O(\log_m n)$，用于数据库索引
\end{enumerate}
\vspace{1em}
\begin{center}
\textcolor{blue}{\textbf{下节预告：排序入门}}
\end{center}
\end{frame}
\end{document}


\subsection{1-18 排序入门}
\documentclass[10pt, aspectratio=169]{beamer}
\usetheme{Madrid}
\usecolortheme{dolphin}
\usepackage{ctex}
\usepackage{xeCJK}
\usepackage{amsmath,amssymb}
\usepackage{listings}
\usepackage{tikz}
\usepackage{xcolor}
\title{408考研零基础 --- 数据结构}
\subtitle{1-18 排序入门}
\author{}
\date{}
\setbeamertemplate{page number in head/foot}{}
\begin{document}
\begin{frame}{本节目标}
\begin{enumerate}
\item 理解排序的基本概念（稳定性、内/外部排序）
\item 掌握插入排序、冒泡排序、简单选择排序
\item 理解快速排序的分治思想
\end{enumerate}
\end{frame}
\begin{frame}{排序的基本概念}
\textbf{稳定性}：相等元素排序后相对顺序不变
\vspace{0.3em}
\textbf{分类}
\begin{itemize}
\item 内部排序：所有元素在内存中
\item 外部排序：元素需借助外存
\end{itemize}
\end{frame}
\begin{frame}{插入排序}
\textbf{思想}：将元素插入已排序区的合适位置
\vspace{0.3em}
\textbf{性能}
\begin{itemize}
\item 最好 $O(n)$（已有序）
\item 最坏 $O(n^2)$（逆序）
\item 平均 $O(n^2)$
\end{itemize}
\vspace{0.3em}
空间 $O(1)$，\textbf{稳定}
\end{frame}
\begin{frame}{冒泡排序}
\textbf{思想}：相邻比较交换，最大元素移至末尾
\vspace{0.3em}
\textbf{性能}
\begin{itemize}
\item 最好 $O(n)$（已有序，提前结束）
\item 最坏 $O(n^2)$（逆序）
\item 平均 $O(n^2)$
\end{itemize}
\vspace{0.3em}
空间 $O(1)$，\textbf{稳定}
\end{frame}
\begin{frame}{简单选择排序}
\textbf{思想}：每趟选最小元素与待排序区第一个交换
\vspace{0.3em}
\textbf{性能}
\begin{itemize}
\item 始终 $O(n^2)$（与初始状态无关）
\item 空间 $O(1)$
\end{itemize}
\textbf{不稳定}
\end{frame}
\begin{frame}{快速排序}
\textbf{思想}：分治 + 分区
\vspace{0.3em}
\textbf{分区操作}
\begin{itemize}
\item 选枢轴（通常第一个元素）
\item low/high 指针交替扫描交换
\item 左边 $\le$ 枢轴 $\le$ 右边
\end{itemize}
\vspace{0.3em}
\textbf{性能}
\begin{itemize}
\item 最好 $O(n\log n)$（均匀分割）
\item 最坏 $O(n^2)$（已有序）
\item 平均 $O(n\log n)$
\end{itemize}
空间 $O(\log n)$ 至 $O(n)$，\textbf{不稳定}
\end{frame}
\begin{frame}{排序算法对比}
\begin{tabular}{|c|c|c|c|c|}
\hline 算法 & 最好 & 平均 & 最坏 & 稳定 \\ \hline
插入排序 & $O(n)$ & $O(n^2)$ & $O(n^2)$ & 稳定 \\ \hline
冒泡排序 & $O(n)$ & $O(n^2)$ & $O(n^2)$ & 稳定 \\ \hline
选择排序 & $O(n^2)$ & $O(n^2)$ & $O(n^2)$ & 不稳 \\ \hline
快速排序 & $O(n\log n)$ & $O(n\log n)$ & $O(n^2)$ & 不稳 \\ \hline
\end{tabular}
\vspace{0.5em}
快速排序是实际应用中最常用的排序算法
\end{frame}
\begin{frame}{本节总结}
\begin{enumerate}
\item 插入排序：插入已排序区，稳定
\item 冒泡排序：相邻交换，稳定
\item 选择排序：每趟选最小，不稳定
\item 快速排序：分治 + 分区，最常用
\item 数据结构模块至此结束，共 18 集
\end{enumerate}
\vspace{1em}
\begin{center}
\textcolor{blue}{\textbf{下个模块：计算机组成原理}}
\end{center}
\end{frame}
\end{document}


\end{document}